% articulo-10-listas.tex - Listas y tablas.
% (enumitem, booktabs, longtable)
%
% Listas ampliadas a 6 niveles con \renewlist + \setlistdepth.
% LaTeX core limita itemize y enumerate a 4; esto lo elimina redefiniendo
% las listas estándar (pandoc genera \begin{itemize} estándar).
%
% Nivel 1: raya sangrada \ParIndent + 2em desde el margen.
%   leftmargin usa \ParIndent directamente (medida fija), NO \parindent,
%   para que funcione correctamente dentro de footnotes donde footmisc
%   y \@parboxrestore pueden alterar \parindent en cualquier momento.
% Niveles 2–6: leftmargin=* (enumitem sitúa la viñeta en la vertical del
%   texto del nivel superior). Viñetas: raya → boliche → ◆, ciclando.
% topsep=partopsep=0pt: elimina el espacio extra de \setstretch antes/después.
%
% \setlistdepth y \renewlist van aquí, en el preámbulo (necesitan ejecutarse
% antes del documento). Los \setlist van en \AtBeginDocument para aplicarse
% después de que todos los paquetes hayan cargado, evitando que alguno los
% sobrescriba.

\usepackage{booktabs}
\usepackage{longtable}
\usepackage{enumitem}

% \tightlist es requerido por pandoc para listas compactas generadas con el
% modo «tight» de Markdown (ítems sin línea en blanco entre ellos).
\providecommand{\tightlist}{%
  \setlength{\itemsep}{0pt}\setlength{\parskip}{0pt}}

\setlistdepth{6}
\renewlist{itemize}{itemize}{6}
\renewlist{enumerate}{enumerate}{6}

\AtBeginDocument{%

\setlist[itemize,1]{
  label      = \textemdash,
  leftmargin = \dimexpr\ParIndent + 1em + 1em\relax,
  labelwidth = 1em,
  labelsep   = 1em,
  itemsep    = 0pt,
  parsep     = 0pt,
  topsep     = 0pt,
  partopsep  = 0pt
}
\setlist[itemize,2]{
  label      = \textbullet,
  leftmargin = *,
  labelwidth = 1em,
  labelsep   = 1em,
  itemsep    = 0pt,
  parsep     = 0pt,
  topsep     = 0pt,
  partopsep  = 0pt
}
\setlist[itemize,3]{
  label      = {\small◆},
  leftmargin = *,
  labelwidth = 1em,
  labelsep   = 1em,
  itemsep    = 0pt,
  parsep     = 0pt,
  topsep     = 0pt,
  partopsep  = 0pt
}
\setlist[itemize,4]{
  label      = \textemdash,
  leftmargin = *,
  labelwidth = 1em,
  labelsep   = 1em,
  itemsep    = 0pt,
  parsep     = 0pt,
  topsep     = 0pt,
  partopsep  = 0pt
}
\setlist[itemize,5]{
  label      = \textbullet,
  leftmargin = *,
  labelwidth = 1em,
  labelsep   = 1em,
  itemsep    = 0pt,
  parsep     = 0pt,
  topsep     = 0pt,
  partopsep  = 0pt
}
\setlist[itemize,6]{
  label      = {\small◆},
  leftmargin = *,
  labelwidth = 1em,
  labelsep   = 1em,
  itemsep    = 0pt,
  parsep     = 0pt,
  topsep     = 0pt,
  partopsep  = 0pt
}

\setlist[enumerate,1]{
  leftmargin = \dimexpr\ParIndent + 1em + 1em\relax,
  labelwidth = 1em,
  labelsep   = 1em,
  itemsep    = 0pt,
  parsep     = 0pt,
  topsep     = 0pt,
  partopsep  = 0pt
}
\setlist[enumerate,2]{
  leftmargin = *,
  labelwidth = 1em,
  labelsep   = 1em,
  itemsep    = 0pt,
  parsep     = 0pt,
  topsep     = 0pt,
  partopsep  = 0pt
}
\setlist[enumerate,3]{
  leftmargin = *,
  labelwidth = 1em,
  labelsep   = 1em,
  itemsep    = 0pt,
  parsep     = 0pt,
  topsep     = 0pt,
  partopsep  = 0pt
}
\setlist[enumerate,4]{
  leftmargin = *,
  labelwidth = 1em,
  labelsep   = 1em,
  itemsep    = 0pt,
  parsep     = 0pt,
  topsep     = 0pt,
  partopsep  = 0pt
}
\setlist[enumerate,5]{
  leftmargin = *,
  labelwidth = 1em,
  labelsep   = 1em,
  itemsep    = 0pt,
  parsep     = 0pt,
  topsep     = 0pt,
  partopsep  = 0pt
}
\setlist[enumerate,6]{
  leftmargin = *,
  labelwidth = 1em,
  labelsep   = 1em,
  itemsep    = 0pt,
  parsep     = 0pt,
  topsep     = 0pt,
  partopsep  = 0pt
}

}% fin \AtBeginDocument
